\chapter{Implementation}\label{ch:implementation}

This chapter defines all implementation choices that were made to validate the research question and test the hypotheses defined in \ref{sec:research-question}.

\section{Tech stack}\label{sec:tech-stack}
\subsection{Hardware}
\subsubsection{Computer}
The development machine was a MacOS laptop with 32 GB of memory and 994 GB of storage. The storage was important to install and run the WebArena configurations, since they ship as Docker images of four full websites and take up a lot of space. This machine ran the private and shared memory generation, against Apple's MLX runtime and the Metal Performance Shaders (MPS) backend. Both models are used here: Gemma for the differentially private token-by-token generation, and Qwen for the title and description that follow it.

\subsubsection{GPUs}
The GPUs that are available to us through the school are a cluster of three NVIDIA A30s. These were used to run the agent, and therefore the Think node, across the 812 WebArena tasks. Although this was not the quickest processing unit available, it was all that was available on a free basis, and each experimental condition requires a complete pass-through of all 812 tasks, so the throughput gain over the laptop compounds across runs. These GPUs were shared across multiple students, and therefore had to be used efficiently.

The work is therefore split across the two machines by what each part needs. The agent loop is throughput-bound, since it repeats across 812 tasks and once per experimental condition, and it runs on the cluster. The differentially private path is bound instead by the token-by-token logit vector prediction over multiple memories, which is very time-consuming, and it runs on the laptop, where the model can be held in memory across a whole generation run without competing for shared cluster time. This split influenced a number of choices on the software side.

\subsection{Software}

\subsubsection{Programming language}
The programming language for this entire project is Python. The reason for choosing it was for multiple reasons: first, it is the programming language the author is most well-versed in. Second, Python hosts a mature scientific stack, built on the array-programming foundation of NumPy \citep{harris2020array} and including the deep learning framework this thesis depends on for local logit access, PyTorch \citep{paszke2019pytorch}. Both the differentially private mechanism and the agent loop are therefore assembled from maintained, interoperable libraries rather than reimplemented.

\subsubsection{Embedding model}\label{sec:embedding-model}
The thesis embeds text with \texttt{all-MiniLM-L6-v2}, a sentence-transformers model \citep{reimers2019sentencebert} built as a six-layer reduction of the distilled MiniLM architecture \citep{wang2020minilm} and fine-tuned with a contrastive objective on roughly 1.17 billion sentence pairs. It is used in two places: to retrieve memories for the Think prompt, and to assign each memory to one of the differentially private labels in the shared memory pipeline described in Section \ref{sec:differential privacy mechanism}.

The model was chosen for its cost rather than its ceiling. At 22.7 million parameters and 384 dimensions it runs offline and deterministically on the laptop, and the retrieval step is called once per Think turn across 812 tasks per pass-through, so the per-call cost compounds quickly. MTEB \citep{muennighoff2023mteb} reports that no single embedding model dominates across task types, so the choice is a trade-off rather than a ranking, and \texttt{all-MiniLM-L6-v2} sits at the small and fast end of it. While a larger encoder would likely retrieve more accurately, this thesis compares configurations against each other under one fixed retrieval setting, so holding retrieval quality constant across conditions matters more than maximising it.

\subsubsection{LLM models}
The LLM models that this thesis uses were chosen with the aim of firstly, meeting the requirements for the use in terms of output, secondly, of being free of charge, thirdly, of producing high-quality output, and fourth, of producing output in a reasonable time-frame. Lastly, this thesis gave preference to models used by either the ReasoningBank or Amin et al. approach, since the aim is to replicate their mechanisms as closely as possible. These objectives had to be traded off against each other for the different use cases.

LLMs were used in the thesis in three parts:
\begin{enumerate}
    \item \textbf{Agent - Thinking node.} The thinking node in the ReAct loop
    requires an LLM to take the information from the past actions and
    observations and make a choice on the next action from a subset of
    pre-defined actions. This step does not require any intermediate
    computations, and simply requires a complete answer: a single line naming
    one action call. The choice of model could therefore be focused on a free
    and quick model.

    In a first step, the author checked for the model that was used to run the
    agent in the ReasoningBank paper. Unfortunately, the paper used
    closed-source, paid models: Gemini-2.5 and Claude 3.7. Since there was no
    budget for this project, and any low-tier free access would have been
    exhausted well before 812 tasks completed, this thesis had to take a
    different approach\footnote{This is an important distinction to the tests
    run by the ReasoningBank paper \citep{ouyang2025reasoningbank}. It would
    therefore not be unexpected if the model underperformed in comparison to
    that paper on the 812 benchmarks.}.

    For open-source models, Qwen is among the most widely adopted open-weight
    model families, with its technical reports accumulating several thousand
    citations \citep{yang2024qwen25, yang2025qwen3} and its checkpoints
    released under Apache 2.0. It could be downloaded for free, and had
    different parameter settings to choose from to balance speed and
    intelligence. The model used for the Think node is Qwen3.5 at 4 billion
    parameters, quantised to NVFP4 (Ollama tag
    \texttt{qwen3.5:4b-nvfp4}).\footnote{Qwen3.5 was released without a
    technical report; the architecture and training lineage are documented for
    the preceding generation in \citet{yang2025qwen3}.}

    The reason for choosing the 4bn parameter model over the 9bn one, is the
    speed of execution. The 9bn parameter model was taking overall too long,
    whereas the 4bn parameter model allowed an average of $\sim$5\,s per Think
    step. Two limits on that figure should be stated. It is a wall-clock
    measurement of one configuration rather than evidence about any single
    factor, because moving to it changed the parameter count and the
    quantisation format at the same time, and no controlled ablation was run.
    It was also measured on the laptop during model selection, whereas the
    evaluation runs were served from the A30 cluster, so it is the basis on
    which the 4bn model was chosen and not the per-step cost of the reported
    runs.

    Since the thesis aim was not to assess overall success rates, but a
    comparative assessment instead, the overall intelligence of the model was
    not as important. Published zero-shot success rates for open-weight models
    at or below 8 billion parameters on the full 812-task benchmark span
    roughly 2\% to 24\%, and that spread is driven at least as much by the
    agent harness, the observation format and the prompting strategy as by the
    model itself \citep{zhang2025symbiotic, gandhi2026gobrowse, lu2026distillation}.
    The closest published reference point for the configuration used here is
    \citet{lu2026distillation}, who report the same Qwen3.5 4bn model at 24.1\%
    on the full 812 tasks under the BrowserGym evaluation protocol, without
    fine-tuning. That figure is a useful target rather than a prediction, since
    the harness differs from the one built for this thesis. While a low
    absolute success rate would narrow the base of tasks on which a difference
    between configurations can show, it would not undermine the comparison,
    which holds the model fixed and varies only the memory condition.

    The model is served locally through an Ollama endpoint, rather than being
    loaded directly into the runner. The reason for this, is that the model is
    quantised, and its quantisation format was built for a newer generation of
    NVIDIA cards than the A30s available here. Loading it directly through
    HuggingFace would only simulate the quantisation on this hardware, so the
    model would keep the accuracy cost of being compressed without running any
    faster.\footnote{The NVFP4 format targets NVIDIA's Blackwell generation.
    HuggingFace's FP-Quant documentation states that on non-Blackwell cards
    pseudoquantisation is mandatory, and that it emulates the effect of the
    quantisation while providing no speed-up.} That would leave a choice
    between a simulated version that is no quicker, and dropping the quantised
    model for the full-size weights. Ollama runs the quantised model as it is
    meant to be run, and avoids the choice altogether.

    Ollama also releases the model from GPU memory once it has been idle for a
    while, whereas a model loaded into the runner holds that memory for as long
    as the process is alive. Since the A30s are shared across students, this
    matters.

    The same model and endpoint serve every other non-private call in the
    system: the success-or-failure judge and the memory extractor described in
    Section \ref{sec:private memory creation experiment}, and the title and
    description generation below. Only the differentially private path departs
    from it, for the reason given next.

    \item \textbf{Differential privacy mechanism.} The model used to generate
    the shared memories could no longer be the Qwen model. The Amin et al.
    mechanism needs the full next-token logit vector over the vocabulary at
    every sampling step, for each member of the batch and for the public
    prompt, and it needs to be able to extend that generation one token at a
    time. Qwen is served in this thesis through Ollama's chat endpoint, which
    returns the generated text and nothing else, so the logit vectors are never
    exposed. The differentially private path consequently runs on a locally
    loaded Gemma model, through PyTorch and HuggingFace transformers, where the
    logits are read directly off the forward pass.

    Naively, each sampled token would cost one full forward pass over the
    prompt for every sensitive prompt plus one for the public prompt. The
    implementation instead stacks all of them into a single padded batch, runs
    one prefill with the key-value cache retained, and extends the generation
    one token at a time through a single-column continuation that reuses that
    cache, so the unchanging prompt prefix is attended to once rather than at
    every step.

    The sole deviation of substance from the \citet{amin2024private} paper is
    that it uses Gemma 1.1 2B IT, whereas this thesis uses the more recent
    Gemma 2 2B IT \citep{gemmateam2024gemma2}. This choice may affect the
    hyperparameters, since the values the paper reports were tuned on the
    earlier model, and the values used here are given in Table
    \ref{tab:dp-hyperparameters} below.

    Gemma 2 2B IT was chosen over the Gemma 2 2B base model for two reasons.
    First, it has been instruction-tuned, meaning further trained on tasks
    phrased as natural-language instructions \citep{wei2022finetuned}, so
    inputs must be arranged into the model-specific chat template. Second, in
    the intermediate experiments reported in Section \ref{intermediate results},
    Gemma 2 2B IT under the chat template produced lessons that fitted the
    input memory items better than the base model under raw completion. That
    comparison was a single qualitative pass over a ten-item fixture, with the
    model and the prompting regime varied together, so it is reported as the
    reason the project standardised on the IT variant rather than as a measured
    result.

    The model is loaded in \texttt{torch.bfloat16} rather than fp32 or fp16.
    Gemma 2 applies logit soft-capping in each attention layer and in the final
    layer \citep{gemmateam2024gemma2}, and fp16's narrower exponent range
    produces silent NaN and Inf values there, while fp32 was too slow on MPS.
    bfloat16 is the model's training dtype and preserves the fp32 dynamic range
    at roughly half the per-forward cost.

    \item \textbf{Shared memory title and description generation.} The Qwen
    model is used once more, for the title and description generation of the
    shared memory. This step no longer requires token logit access, so it can
    be served through the same Ollama endpoint as the Think node. Its inputs
    are strictly the differentially private content and the differentially
    private label, so by the post-processing property of differential privacy
    the call adds nothing to the privacy account.

    The choice of Qwen over Gemma at this step is driven by role rather than by
    a measured capability gap. Gemma is loaded for the differentially private
    path because that path needs token-level logit access, which is a
    requirement of the mechanism and not a judgement about the model. Once the
    content has been released, the remaining work is short-form instruction
    following, and Qwen is already resident in the pipeline behind the same
    \texttt{ChatClient} interface the Think node uses, so reusing it costs
    nothing. No head-to-head benchmark between the two is reported here: a 4
    billion parameter model released in 2026 and a 2 billion parameter model
    released in mid-2024 differ in scale, training recipe and release date at
    once, so a comparison between them would be confounded on three axes and
    its outcome is not in question.
\end{enumerate}

\subsubsection{Testing configuration}
The testing configuration consisted of both the tasks and scoring, as well as the environment. The tasks and scoring were provided by WebArena, while the environment was set by BrowserGym.

\textbf{WebArena}

All evaluation runs are completed on the WebArena gym \citep{zhou2024webarena},
accessed through the BrowserGym wrapper \citep{dechezelles2024browsergym}. The
WebArena repository provides the evaluation environment, tasks, and scoring for
812 different online tasks:

\begin{itemize}
    \item \textbf{Environment.} The WebArena repository provides the
    configurations for four websites, which mimic the real-world equivalent of
    GitLab, Reddit, an online store content management system (CMS), and
    OneStopShop. These configurations allow the tester to run the agent tests
    locally on replicas of the websites, removing the risk that the agent's
    actions take effect on a live service. It also provides a set of utility
    tools and knowledge resources that the agent can call, if needed, for
    further closeness to the real world scenario. These include a calculator, a
    scratchpad, an offline copy of Wikipedia, and a map, the last of which is
    itself served as a site and carries its own tasks \citep{zhou2024webarena}.
    Because the environment resets to a deterministic initial state, these
    configurations also ensure that the same tests can be run over an extended
    duration of time, without fear of change to the website.

    \item \textbf{Tasks.} The repository also provides a list of 812 different
    online tasks, which the agent is given to perform. These tasks are
    instantiated from 241 templates and are phrased as high-level, realistic
    natural language commands rather than click-by-click instructions, so the
    agent has to decompose them itself \citep{zhou2024webarena}. The example
    task \textit{``What is the top-1 best-selling brand in Quarter 1 2022''}
    names an outcome and leaves the route to it open.

    \item \textbf{Evaluation.} Lastly, the repository also provides a program,
    which validates the functional correctness of each task by checking the
    outcome rather than the action sequence taken to reach it. BrowserGym
    surfaces that verdict as the episode reward, and the runner records it when
    the episode ends, so a task counts as successful when its final reward is
    greater than zero.
\end{itemize}

The reason for choosing the WebArena environment was to keep as close as possible to the original testing environment used by \citet{ouyang2025reasoningbank}. The thesis therefore chose to test its hypotheses on the same test set as that paper. Firstly, this ensures that the ReasoningBank framework is useful on this dataset, and that a difference in agent behaviour with memories should be observable. Second, it allows the results of this thesis to be compared to those in the paper.

\textbf{BrowserGym}

While WebArena provides the benchmark tasks and evaluations, BrowserGym provides the testing environment to run the thesis on. It launches and takes action on a Chromium browser using Playwright. While the work for this thesis included creating its own Playwright nodes and Chromium actions, BrowserGym simplifies the work in three ways. First, it exposes the WebArena suite through a single environment call, including the per-task setup and the reward signal that the hand-rolled path had no access to. Second, it offers a substantially larger action vocabulary than the seven actions the hand-rolled implementation supported, set out in Section \ref{sec:experiment_design}. Third, it catches exceptions raised by an action and returns the error text in the next observation \citep{dechezelles2024browsergym}, so the failure is fed back to the LLM rather than lost.

It also returns structured observations, in one of three formats. BrowserGym automatically returns all three observation spaces which WebArena has been tailored to \citep{zhou2024webarena, dechezelles2024browsergym}: the HTML DOM tree, a screenshot, or the accessibility tree of the website URL and content. Which of the three this thesis feeds into the Think node, and the reason for that choice, is set out in Section \ref{sec:experiment_design}.

%%%%%%%%%%%%
\section{Design choices}\label{sec:experiment_design}
This section lays out how the system architecture and algorithm from the previous sections were implemented to assess this thesis' research question. While Section \ref{sec:tech-stack} sets out which tools and models the thesis uses, this section sets out how they were configured.

The hypotheses listed in Section \ref{sec:hypotheses} require the system architecture described above to be run on the 812 WebArena tasks\footnote{This thesis will refer to a pass-through as a complete evaluation of all 812 WebArena tasks under one configuration.} with different configurations. To assess the benefit of shared memories, the thesis requires a pass-through of the following configurations:
\begin{itemize}
    \item \textbf{A - No memory use.} The agent runs without private or shared memories. The trajectories from this run are used to generate private and shared memories for future pass-throughs.

    \item \textbf{B - Private memory use only.} In this pass-through all memories from the previous 812 trajectories are considered to have been created under a single user. This next pass-through is considered to be done from the same user, who therefore has access to all memories from the past pass-through, including the memory from the exact trajectory's pass-through. The agent thus has access to all private memories using a top-3 retrieval.

    Because a memory is retrieved on the embedding of the task it was distilled from, and that task is the identical string to the one being solved, a task's own memory ranks first on every retrieval. This pass-through therefore measures the ceiling of what memory injection can contribute, and not the transfer between different tasks.
    % TODO(Yasmine): confirm --k 3; the runner's CLI default is 1.

    \item \textbf{C - Shared memories only.} Using the private memories generated under pass-through A, the shared memory mechanism generates the new shared memories. The agent has access to the shared memories only, and none of the private memories for the 812 tasks.
    % TODO(Yasmine): state how many shared memories the cycle produced. With
    % three labels and one memory per label, the ceiling is three per cycle.

    \item \textbf{D - Private and shared memories.} This last pass-through is used to compare how private and shared memories are used when both are available, and whether the shared memories are helpful or detrimental to the overall functioning of the system. This is a complete implementation of the architecture described in section \ref{sec:system architecture}.
\end{itemize}

Each of these pass-throughs requires the agent component of the system. The private and shared memory creation components run once, on the trajectories produced by pass-through A, and their outputs are what pass-throughs B, C and D consume. Next, this thesis describes the implementation choices for each step.

\subsection{Agent}
While WebArena provides the set-up, the ReAct Loop is the engine underlying the agent completing the tasks. The implementation of these is described in more detail below.

\subsubsection{ReAct Loop}
The ReAct Loop is run in line with the academic description provided in
\ref{sec:react-framework}. The agent loops through observation-thinking-action steps in order to solve the task it has been given:
\begin{itemize}
    \item The \textbf{observation} step is handled by BrowserGym and the
    WebArena configuration. Although BrowserGym returns an accessibility tree, the HTML DOM, and a screenshot of the website url and content, this thesis only feeds the accessibility tree observation into the LLM's thought process.
    The reason for choosing the accessibility tree over the HTML DOM or
    screenshot representation is that it contains the structure of the page at
    the lowest cost to the context window, and in a form the action grammar can consume directly. The accessibility tree is a subset of the DOM restricted to the elements that are relevant and useful for displaying a page's content, so it retains that structure at a fraction of the DOM's length \citep{zhou2024webarena}. Its compactness is highly relevant for the class of model used here, since the thesis uses lower-capability models, and \citet{enomoto2026observation} find on WorkArena that accessibility-tree observations raise the success rate of lower-capability
    open-source models. In comparison, HTML only helps the higher-capability proprietary ones, because the longer input drives up hallucination in the weaker models.
    The tree also labels each element with an identifier that the BrowserGym
    action grammar takes directly, so the LLM can name a click target without
    grounding it in pixel coordinates, which a screenshot would require of a
    vision model this thesis does not use.
    A final implementation choice that this thesis made, is to cap the rendered tree at 12,000 characters, since even the accessibility tree can outgrow the context window on a dense page. The truncation is undertaken from the end and marked with an explicit truncation comment, so that page content never displaces the system prompt or the task.

    \item The \textbf{think} process is handled by the LLM.
    Each call is made at temperature 0, with the model's reasoning trace
    disabled and the reply capped at 64 tokens, which is the length of a single
    action line. Holding the sampling settings
    fixed keeps the Think step deterministic given its prompt, so any change in
    success rate between evaluation runs is attributable to the memories rather
    than to sampling noise.

    The reasoning trace is disabled because the reply grammar has no room for
    it. Qwen emits its reasoning as a block before its answer, whereas the
    system prompt asks for exactly one line holding a single function call, so a
    trace would either have to be stripped before parsing or would be rejected
    outright and spend a step. The two settings are coupled: a reasoning block
    would exhaust the 64-token budget before the model reached the action.
    Disabling it also bounds the cost of a pass-through, since the Think step is
    called up to 30 times per task, across 812 tasks, once per configuration.
    While this keeps each pass-through within the time available on a shared
    cluster, it also removes the surface on which a retrieved memory could act
    before an action is chosen, which is a limitation this thesis returns to in
    Chapter \ref{sec:discussion}.
    % PARKED FOR CHAPTER 6 (Discussion): the sensitivity limitation named above.
    % If memories help mainly by shaping reasoning, disabling the reasoning
    % trace may suppress the effect this thesis sets out to measure. Write this
    % up as a named threat to sensitivity once the results are known.

    The Think module is provided the following information, as part of the system
and user prompt:
\begin{itemize}
    \item \textbf{System Prompt.} The initial set of instructions for each
    thinking step is consistent across each agent loop, and consists of three
    separate parts.
    \begin{itemize}
        \item Firstly, it informs the LLM that it is a web agent controlling a
        browser through BrowserGym WebArena actions.
        \item Second, it fixes the reply format, which is exactly one line
        holding a single Python function call and no commentary, with element
        identifiers taken from the accessibility tree.
        \item Third, it lists all the possible actions (see table
        \ref{tab:webarena_action_space}) that the LLM can return, with one
        worked example each. This list is generated directly from the BrowserGym
        action subset rather than written by hand.
    \end{itemize}
    \item \textbf{User Prompt.} The user prompt passes all other instructions
    that are relevant to the agent. These define the problem, rather than the
    LLM's own system.
    \begin{itemize}
        \item Task. The first item in the user prompt states the agent's task, which is a copy of the task provided by WebArena, read verbatim when the
        environment is reset.
        \item Memory. Not all evaluation runs include memories. Specifically,
        the baseline evaluation does not include any memories, so that the
        thesis has a clear understanding of how successful the agent can be
        based on the defined ReAct framework, the prompts, the libraries and the
        models. Later evaluations then append the memories. These are the key
        addition to the evaluation and should provide helpful context to the LLM
        for improving its functioning. They are appended as individual line
        items with both the memory's title and content, prefixed by the
        memory-injection instruction reproduced verbatim from
        \citet{ouyang2025reasoningbank} so that prompt wording is held fixed as
        a control variable.
        \item Agent trajectory. The agent also receives a history of past steps,
        thoughts, and outcomes from this trajectory. The outcome of each step is
        either \texttt{ok}, if the step was completed without an error message,
        or the specific failure, which is the error returned by BrowserGym, the
        exception the step raised, or the parse error the reply produced. These
        are the working memory which enable it to learn from previous actions,
        and work towards its goal.
        \item Accessibility tree. As discussed above, the LLM is passed the
        website observation in the form of an accessibility tree.
        \item Instruction. The end of the user prompt is a single sentence:
        \textit{``Reply with exactly one BrowserGym action line.''} This pushes
        the LLM to respond with a single action item that can be parsed for the
        action step.
    \end{itemize}
\end{itemize}
    \item The \textbf{action} space is defined by WebArena once more, and
    includes actions across three categories. The first includes element
    operations including clicking, hovering, and typing, while the second
    enables tab-related actions. Lastly, the third allows URL navigation
    actions. The action space as published is listed in table
    \ref{tab:webarena_action_space} below.
    The set actually exposed to the LLM is BrowserGym's \texttt{webarena}
    subset, which departs from the published table in three places. It adds
    \texttt{select\_option} for dropdown elements. It replaces the paper's
    \texttt{stop(answer)} with two explicit termination actions,
    \texttt{send\_msg\_to\_user} for returning an answer and
    \texttt{report\_infeasible} for declaring the task unachievable. And it
    omits \texttt{noop}. Both the grammar printed into the system prompt and the
    parser that validates the LLM's reply are generated from this one subset, so
    the actions the LLM is offered and the actions the environment accepts
    cannot drift apart.
\end{itemize}

The agent continues looping through the ReAct framework until the environment
terminates, which happens when the LLM emits \texttt{send\_msg\_to\_user} or
\texttt{report\_infeasible}. Meeting the goal and declaring the task infeasible
are both specified as possible actions for the LLM to take as part of its Think
process, so ending the episode is a choice the model makes rather than a state
the harness infers.

In addition, to prevent continuous looping, this thesis' implementation also
added a stopping rule after 5 successive identical replies from the LLM, and a
max 30 step rule for each task. Five identical replies is the shape that both
observed failure modes take, whether a reply the parser keeps rejecting or a
click repeated on a page whose state never advances. A parse failure on its own
does not end the episode, since the rejected reply and its error are appended to
the history and the loop returns to the observation step, so the LLM sees its
own malformed output on the next turn and can correct it.
While these rules may reduce the agent's overall capability, they bound the time
necessary for each evaluation of all 812 tasks, which has to be repeated once
per experimental condition.

\begin{table}[t]
\centering
\begin{tabular}{l|l}
\hline
\textbf{Action Type} & \textbf{Description} \\
\hline
\texttt{noop} & Do nothing \\
\texttt{click(elem)} & Click at an element \\
\texttt{hover(elem)} & Hover on an element \\
\texttt{type(elem, text)} & Type to an element \\
\texttt{press(key\_comb)} & Press a key comb \\
\texttt{scroll(dir)} & Scroll up and down \\
\hline
\texttt{tab\_focus(index)} & focus on $i$-th tab \\
\texttt{new\_tab} & Open a new tab \\
\texttt{tab\_close} & Close current tab \\
\hline
\texttt{go\_back} & Visit the last URL \\
\texttt{go\_forward} & Undo \texttt{go\_back} \\
\texttt{goto(URL)} & Go to URL \\
\hline
\end{tabular}
\caption{Action Space of \texttt{WebArena}}
\label{tab:webarena_action_space}
\end{table}

Lastly, each completed task writes one row recording the task intent, the full
trajectory, the final observation, the agent's returned answer and the harness
reward. Counting the rows whose reward is greater than zero gives the number of successful and failed trajectories for that evaluation run.

\subsection{Private memory creation}\label{sec:private memory creation experiment}
As described above, the private memory creation mechanism first assesses whether a trajectory was successful or failed, and only in a second step creates the memories from this trajectory.

\subsubsection{Success-or-failure assessment}
The failure-or-success assessment is an LLM-as-judge call, and is separate from the reward check that WebArena performs. Both verdicts are recorded for every trajectory, but the WebArena assessment is what the evaluation in chapter \ref{sec:results} scores. The judge is based on the following prompt from the ReasoningBank paper, shown in table \ref{tab:agent-eval-prompt}.

\begin{table}[htbp]
    \centering
    \caption{System instruction and input prompt template used by the LLM judge to evaluate web-navigation agent trajectories.}
    \label{tab:agent-eval-prompt}
    \small
    \begin{tabularx}{\textwidth}{@{} X @{}}
        \toprule
        \multicolumn{1}{@{}c@{}}{\textbf{System Instruction}} \\
        \midrule
        You are an expert in evaluating the performance of a web navigation agent. The agent is designed to help a human user navigate a website to complete a task. Given the user's intent, the agent's action history, the final state of the webpage, and the agent's response to the user, your goal is to decide whether the agent's execution is successful or not.
        \newline\newline
        There are three types of tasks:
        \begin{enumerate}
            \item Information seeking: The user wants to obtain certain information from the webpage, such as the information of a product, reviews, map info, comparison of map routes, etc. The bot's response must contain the information the user wants, or explicitly state that the information is not available. Otherwise, e.g. the bot encounters an exception and respond with the error content, the task is considered a failure. Besides, be careful about the sufficiency of the agent's actions. For example, when asked to list the top-searched items in a shop, the agent should order the items by the number of searches, and then return the top items. If the ordering action is missing, the task is likely to fail.
            \item Site navigation: The user wants to navigate to a specific page. Carefully examine the bot's action history and the final state of the webpage to determine whether the bot successfully completes the task. No need to consider the bot's response.
            \item Content modification: The user wants to modify the content of a webpage or configuration. Carefully examine the bot's action history and the final state of the webpage to determine whether the bot successfully completes the task. No need to consider the bot's response.
        \end{enumerate}
        \vspace{0.4em}
        *IMPORTANT*
        \newline
        Format your response into two lines as shown below:
        \newline
        Thoughts: \textless your thoughts and reasoning process\textgreater{}"
        \newline
        Status: "success" or "failure"
        \\
        \midrule
        \multicolumn{1}{@{}c@{}}{\textbf{Input Prompt}} \\
        \midrule
        \textbf{User Intent:} \{intent\}
        \newline\newline
        \textbf{Trajectory:} \{trajectory\}
        \newline\newline
        \textbf{The detailed final state of the webpage:}
        \texttt{\textasciigrave\textasciigrave\textasciigrave md \{cap\} \textasciigrave\textasciigrave\textasciigrave}
        \newline\newline
        \textbf{Bot response to the user:} \{response if response else "N/A"\}
        \\
        \bottomrule
    \end{tabularx}
\end{table}

The four insertions into this system prompt are the user intent, the trajectory, the final state, and the bot response:
\begin{itemize}
    \item The user intent - i.e. the task description -, and the trajectory are both inserted into the judge's system prompt from the trajectory record.
    \item The third addition into the prompt is the final state of the webpage, which isn't part of the trajectory record. The final state is inserted as an accessibility tree of the webpage as the agent last saw it, truncated to the same 12,000-character budget.
    \item Lastly, the bot response is the LLM's last action. If the action was the \texttt{send\_msg\_to\_user} action, then the message itself is attached as the response, whereas any other action would result in the literal \textit{``N/A''} as the bot response.
\end{itemize}

One failure mode of the extraction is handled explicitly. A judge reply that carries no parseable \textit{``Status''} line is treated as a failure, so a trajectory whose success the LLM did not confirm is never routed to the success prompt.

\subsubsection{Memory creation}
After marking a trajectory as either successful or failed, a separate instantiation of the LLM  then creates up to three summary memories using the prompts set out in the figure below \ref{fig:memory-extraction-prompts}. Both prompts are reproduced verbatim from \citet{ouyang2025reasoningbank}.
An extraction that yields no complete memory item, either because the reply carried no item header or because every item was missing one of the three fields, causes the trajectory to be dropped rather than stored as an empty record.

The extractor call runs at temperature 1.0, reproduced from \citet{ouyang2025reasoningbank}, rather than at the temperature 0 used everywhere else. The memories are therefore drawn stochastically, so a second extraction pass over the same trajectories would not reproduce them exactly. Holding the setting fixed to the published value keeps the extraction comparable to ReasoningBank, at the cost of making the memory set itself a sampled artefact rather than a deterministic function of pass-through A.

\definecolor{succgreen}{RGB}{0,128,0}
\definecolor{failred}{RGB}{200,30,30}
\newcommand{\hlsuccess}[1]{\textcolor{succgreen}{\textbf{#1}}}
\newcommand{\hlfail}[1]{\textcolor{failred}{\textbf{#1}}}
% Monochrome alternative:
% \renewcommand{\hlsuccess}[1]{\textbf{#1}}
% \renewcommand{\hlfail}[1]{\textbf{#1}}
\newcommand{\bt}{\texttt{\textasciigrave\textasciigrave\textasciigrave}}
\begin{figure}[htbp]
    \centering
    \footnotesize
    \noindent%
    %-----------------------------------------------------------------
    % LEFT PANEL: successful trajectories
    %-----------------------------------------------------------------
    \begin{minipage}[t]{0.475\textwidth}%
        \begin{tabularx}{\linewidth}{@{} X @{}}
            \toprule
            \multicolumn{1}{@{}c@{}}{\textbf{System Instruction}} \\
            \midrule
            You are an expert in web navigation. You will be given a user query, the corresponding trajectory that represents \textbf{how an agent} \hlsuccess{successfully accomplished the task}.
            \newline\newline
            \#\# Guidelines
            \newline
            You need to extract and summarize useful insights in the format of memory items \hlsuccess{based on the agent's successful trajectory}.
            \newline
            The goal of summarized memory items is to be helpful and generalizable for future similar tasks.
            \newline\newline
            \#\# Important notes
            \begin{itemize}[label=--, leftmargin=1.2em, itemsep=0.15em, topsep=0.2em]
                \item You must first \hlsuccess{think why the trajectory is successful}, and then summarize the insights.
                \item \hlsuccess{You have learned or seen something new from the successful trajectory.}
                \item You can extract \textit{at most 3} memory items from the trajectory.
                \item You must not repeat similar or overlapping items.
                \item Do not mention specific websites, queries, or string contents, but rather focus on the generalizable insights.
            \end{itemize}
            \vspace{0.3em}
            \#\# Output Format
            \newline
            Your output must strictly follow the Markdown format shown below:
            \newline
            \bt
            \newline\newline
            \# Memory Item i
            \newline
            \#\# Title \textless the title of the memory item\textgreater{}
            \newline
            \#\# Description \textless one sentence summary of the memory item\textgreater{}
            \newline
            \#\# Content \textless 1-3 sentences describing the insights learned to successfully accomplishing the task\textgreater{} \bt
            \\
            \midrule
            \multicolumn{1}{@{}c@{}}{\textbf{Input Prompt}} \\
            \midrule
            \textbf{Query:} \textless user query\textgreater{}
            \newline\newline
            \textbf{Trajectory:} \textless trajectory that completes the query\textgreater{}
            \\
            \bottomrule
        \end{tabularx}%
    \end{minipage}%
    \hfill%
    %-----------------------------------------------------------------
    % RIGHT PANEL: failed trajectories
    %-----------------------------------------------------------------
    \begin{minipage}[t]{0.475\textwidth}%
        \begin{tabularx}{\linewidth}{@{} X @{}}
            \toprule
            \multicolumn{1}{@{}c@{}}{\textbf{System Instruction}} \\
            \midrule
            You are an expert in web navigation. You will be given a user query, the corresponding trajectory that represents \textbf{how an agent} \hlfail{attempted to resolve the task but failed}.
            \newline\newline
            \#\# Guidelines
            \newline
            You need to extract and summarize useful insights in the format of memory items \hlfail{based on the agent's failed trajectory}.
            \newline
            The goal of summarized memory items is to be helpful and generalizable for future similar tasks.
            \newline\newline
            \#\# Important notes
            \begin{itemize}[label=--, leftmargin=1.2em, itemsep=0.15em, topsep=0.2em]
                \item You must \hlfail{first reflect and think why the trajectory failed}, and then summarize what lessons you have learned or strategies to prevent the failure in the future.
                \item You can extract \textit{at most 3} memory items.
                \item You must not repeat similar or overlapping items.
                \item Do not mention specific websites, queries, or string contents, but rather focus on the generalizable insights.
            \end{itemize}
            \vspace{0.3em}
            \#\# Output Format
            \newline
            Your output must strictly follow the Markdown format shown below:
            \newline
            \bt
            \newline\newline
            \# Memory Item i
            \newline
            \#\# Title \textless the title of the memory item\textgreater{}
            \newline
            \#\# Description \textless one sentence summary of the memory item\textgreater{}
            \newline
            \#\# Content \textless 1-3 sentences describing the insights learned to successfully accomplishing the task\textgreater{} \bt
            \\
            \midrule
            \multicolumn{1}{@{}c@{}}{\textbf{Input Prompt}} \\
            \midrule
            \textbf{Query:} \textless user query\textgreater{}
            \newline\newline
            \textbf{Trajectory:} \textless trajectory that completes the query\textgreater{}
            \\
            \bottomrule
        \end{tabularx}%
    \end{minipage}%
    \caption{System instructions for extracting memory items from agent trajectories: the left panel targets successful trajectories (summarizing why they succeed), while the right targets failed trajectories (reflecting on failure and deriving lessons).}
    \label{fig:memory-extraction-prompts}
\end{figure}

The three memories returned by the LLM, which each include a title, description and content, are bundled together with the user ID, the success/failure judgement, and the initial task. Each task is then embedded with the sentence transformer described in section \ref{sec:embedding-model}, producing a 384-dimensional dense vector, and added to the bundle. The embedding is computed once, when the record is written, and stored inside the record, so reloading the store never re-embeds. At retrieval time only the incoming task is embedded, and it is compared against every stored vector by cosine similarity. As stated previously, the aim is to ensure the private memory creation resembles the ReasoningBank architecture as closely as possible for comparability, and ReasoningBank likewise embeds the task string and retrieves by cosine search.

\subsection{Shared memory creation}
The implementation questions around the shared memory creation regarded the privacy hyperparameter choice, the privacy budget, the parameters in the four steps of the shared memory creation, and the prompt choice.

\subparagraph{Hyperparameters.} The privacy hyperparameters are held identical across Step 1 and Step 3, so that the two rounds are directly comparable and their budgets add without further analysis. The values in Table \ref{tab:dp-hyperparameters} are the settings of the pilot configuration and are for a theoretical, rather than practical proof-of-concept. The $\varepsilon$ in particular is set far above any privacy-meaningful level.

% INTERIM TABLE
\begin{table}[t]
\centering
\begin{tabular}{l|l|l}
\hline
\textbf{Symbol} & \textbf{Meaning} & \textbf{Pilot value} \\
\hline
$c$ & Logit clipping bound & 50.0 \\
$\tau$ & Sampling temperature, private branch & 1.0 \\
$\tau_{\text{public}}$ & Sampling temperature, public branch & 1.5 \\
$\sigma$ & Laplace noise scale for the threshold & 0.1 \\
$\theta$ & Public/private distance threshold & 0.0 \\
$r_{\max}$ & Cap on private tokens per output & 80 \\
$\varepsilon$ & Privacy budget per round & 200.0 \\
$\delta$ & Privacy $\delta$ per round & $1/s$ \\
\hline
\end{tabular}
\caption{\textbf{Interim.} Differential privacy hyperparameters of the pilot configuration, applied identically to Step 1 and Step 3. These values are provisional and are superseded once the budget is retuned.}
\label{tab:dp-hyperparameters}
\end{table}

\subparagraph{Privacy budget.} The privacy budget is not spent directly. For a given target $\varepsilon$ and batch size $s$, the code solves for the number of private tokens $r$ that fit inside that budget. Generation then stops once $r$ private tokens have been sampled. Because $\varepsilon$ is accounted from the budgeted $r$ rather than from the tokens the generation actually consumed, a memory that terminates early is reported at the full cost of its budget. The reported $\varepsilon$ is therefore an upper bound on what was spent. Two further conditions halt generation before $r$ is exhausted: a total cap of 256 tokens across the private and public branches combined, since this seems sufficiently long for a memory and caps the generation duration, and the model emitting either its end-of-sequence token or the end-of-turn token of the instruction-tuned chat template.

\subparagraph{Step 1 - Labels per cycle.}
The number of labels generated per cycle is 80. The choice of creating 80 labels is a simple deduction that for a batch size of 30, where each task yields 3 memories to a batch, one would need approximately 10 tasks per label. The batch size of 30 is what then determines the privacy budget.

\subparagraph{Step 1 - Prompt.}
The prompt in Table \ref{tab:label-prompt} is the prompt used in the differentially private mechanism in Step 1 to generate the labels.

\subparagraph{Step 1 - Output parsing.} The output of the mechanism is a JSON array of label strings, which is parsed with three successive attempts. The output is first parsed as JSON directly. Failing that, a closing bracket is appended, since differentially private sampling can cut the output short before the model reaches it. If this still fails, every double-quoted substring in the output is taken as a label candidate. Any slot that none of the three attempts fills is padded with a placeholder label, and flagged as a fallback in the audit record. A placeholder label will rarely attract memories in Step 2, and the memories that would have been assigned to it are carried into the next cycle rather than lost.

\subparagraph{Step 2 - Minimum number of memories for generation.}
The threshold x, the minimum number of memories a label must attract before a shared memory is generated from it, is set to 5. Every memory in a bucket that does not reach five is written to a carry-over file and re-enters the mechanism at the next cycle.

\subparagraph{Step 3 - Number of memories generated per label.}
The mechanism produces exactly one shared memory per qualifying label per cycle. While the Amin et al. paper runs Algorithm 1 repeatedly to build up a dataset of many synthetic examples from the same private corpus, this thesis focuses on the simplest implementation with one memory. The number of shared memories available to the agent grows across cycles instead, as new private memories accumulate and new labels are synthesised.

\subparagraph{Step 4 - Title and description generation.}
The title and description are generated in a single Qwen call at temperature 0. The reply is rejected if the title is absent, runs beyond ten words, or ends in punctuation, and if the description is absent. On rejection the call is repeated once, and if the second reply also fails validation the title and description are derived from the content itself: the title becomes its first eight words, and the description its first sentence. The fallback is deterministic and takes its input from the content released in Step 3, so it is a post-processing step like the LLM call it replaces and costs no additional budget.

\subparagraph{Retrieval of the shared memories.}
Once a cycle has written its shared memories, they become available to the agent alongside its own. The final evaluation run over the 812 tasks retrieves from both stores, so that each agent sees its own private memories and the shared memories the mechanism produced. A second run retrieves from the shared store alone. That run is diagnostic rather than a condition of either hypothesis: it separates a gain that comes from the shared tier from one that comes from the agent's own history, which the combined run on its own cannot distinguish.

\subparagraph{System prompts.} All system prompts for the shared memory creation are included in the tables below.

% Step 1 system prompt
\begin{table}[htbp]
    \centering
    \caption{Step 1 - System prompt used by the Gemma model to generate the differentially private labels for each cycle.}
    \label{tab:label-prompt}
    \small
    \begin{tabularx}{\textwidth}{@{} X @{}}
        \toprule
        \multicolumn{1}{@{}c@{}}{\textbf{System Prompt}} \\
        \midrule
        Return exactly {k} labels that summarise the memory items below. Each label must be 1–4 words and on a DISTINCT facet — no two labels may be synonyms, paraphrases, or describe the same aspect. Return only a JSON array of strings.
        No explanations, numbering, markdown, or extra text.\newline
        \newline
        Example:\newline
        ["pricing risk", "customer churn", "supply delay", '
        '"budget pressure", "staff training", "data quality"]\newline
        \newline
        Memory items:\newline
        \{items\}
        \bottomrule
    \end{tabularx}
\end{table}

% Step 3 system prompt
\begin{table}[htbp]
    \centering
    \caption{Step 3 - System prompt used by the Gemma model to generate the differentially private content for each label bucket.}
    \label{tab:content-generation-prompt}
    \small
    \begin{tabularx}{\textwidth}{@{} X @{}}
        \toprule
        \multicolumn{1}{@{}c@{}}{\textbf{System Prompt}} \\
        \midrule
        You will be given memory items distilled from one web-navigation
        trajectory related to {label}. Produce a single short paragraph
        (1 to 3 sentences) that captures the most generalisable lesson
        from these memory items. Output only the lesson text — no title,
        no header, no markdown formatting, no list, no explanation of
        what you are doing.\newline
        \newline\newline
        Memory items:\newline
        \{items\}
        \newline\newline
        Lesson:
        \bottomrule
    \end{tabularx}
\end{table}

% Step 4 system prompt
\begin{table}[htbp]
    \centering
    \caption{Step 4 - System prompt used by the Qwen model to generate a title and description for each synthesised shared-memory item.}
    \label{tab:titling-and-description-prompt}
    \small
    \begin{tabularx}{\textwidth}{@{} X @{}}
        \toprule
        \multicolumn{1}{@{}c@{}}{\textbf{System Prompt}} \\
        \midrule
        You will be given the content field of a shared memory item synthesised from one web-navigation topic (label: \{label\}). Write a title (1 to 10 words, capitalised, no trailing punctuation) and a one-sentence description that together identify what the content advises. Do not introduce any facts not present in the content. Output exactly two lines in this format and nothing else:
        \newline\newline
        Title: \textless title\textgreater{}
        \newline
        Description: \textless description\textgreater{}
        \\
        \bottomrule
    \end{tabularx}
\end{table}
